\chapter{Conclusão} \label{Cap:Conclusao}

Este trabalho demonstrou que é possível construir modelos preditivos de churn com qualidade discriminativa relevante a partir de variáveis comportamentais simples, calculadas sobre dados transacionais reais de e-commerce brasileiro. O modelo selecionado atingiu ROC-AUC de 0,8106 na partição de referência, valor que deve ser lido em conjunto com duas estimativas convergentes de desempenho esperado: a validação cruzada de cinco folds sobre o conjunto de treino, conduzida com reamostragem interna a cada partição, resultou em 0,6814, e a média de trinta partições estratificadas independentes, em 0,6686 ($\pm$ 0,0831). A partição de referência situa-se, portanto, na faixa superior da distribuição amostral, e o desempenho esperado em uso é o indicado pelas duas últimas estimativas. O pipeline proposto, com separação rigorosa de janelas temporais e estratégia combinada de balanceamento, mostrou-se eficaz para lidar com os principais desafios metodológicos do problema: data leakage, desbalanceamento extremo de classes e viés de avaliação.

Cabe registrar que a escolha do Random Forest não se apoiou na maximização de métricas de classificação binária. Sob as três métricas recomendadas por De la Cruz Huayanay, Bazán e Russo (\citeyear{delacruz2025}), o modelo indicado seria o HistGradient Boosting, e a liderança isolada em MCC e Kappa de Cohen cabe à Regressão Logística. A adoção do Random Forest fundamenta-se na natureza do entregável: a operação de retenção proposta requer uma priorização de clientes em faixas de risco, e não uma decisão binária, o que desloca o critério para a qualidade da ordenação dos scores, medida por ROC-AUC e pela Average Precision da classe Churn, nas quais o Random Forest lidera. A qualidade dessa ordenação sustenta a estratificação em quatro faixas interpretáveis, cuja validade é aferida pela taxa de churn observada em cada uma, crescente de 22,2\% a 99,9\% na base completa e preservada em direção no conjunto de teste. Verificada fora da amostra, a estratificação retém nove dos dez clientes Ativos do conjunto de teste fora da faixa Crítico, em 39,0\% das observações, lift de 2,31 vezes, magnitude que se mantém quando a mesma métrica é calculada sobre a base completa (2,34 vezes). Em cobertura fixa, critério que neutraliza o tamanho do recorte e isola a ordenação, o Random Forest lidera ou empata os dois concorrentes em todas as frações avaliadas entre 10\% e 60\% do conjunto de teste, e alcança 90\% dos clientes Ativos abordando 35,6\% das observações, contra 49,6\% do HistGradient Boosting e 52,5\% da Regressão Logística. Cabe registrar que recortes mais restritivos, como o das duas faixas de menor risco, apresentam na base completa valores substancialmente superiores aos verificados no conjunto de teste, razão pela qual não são adotados como evidência. O HistGradient Boosting, em contraste, concentra os scores junto do extremo superior da escala (mediana de 0,9621), alocando 77,6\% da base na faixa Crítico, e não expõe medida nativa de importância de features, sobre a qual se apoiam a hierarquia de drivers e a evidência de estabilidade entre partições apresentadas na Seção \ref{sec:interpretabilidade}. A Regressão Logística, embora produza a distribuição de scores mais uniforme entre as quatro faixas, é a menos eficaz das três na identificação de clientes recuperáveis por unidade de cobertura, o que evidencia que a dispersão da escala não constitui, isoladamente, critério de escolha. Registra-se, por fim, que o HistGradient Boosting apresenta ROC-AUC média superior em validação cruzada (0,7363 contra 0,6814) e menor diferença em relação ao hold-out (0,0252 contra 0,1293), conforme a Tabela \ref{tab:benchmark}, diferença que não altera a ordenação em cobertura fixa, critério em que se apoia a seleção.

A análise de correlação de Spearman entre os rankings dos dez modelos revelou comportamento heterogêneo entre as métricas: o ranking por ROC-AUC apresenta correlação moderada com o de Kappa de Cohen (rho = 0,6485, única significativa ao nível de 5\%, com p = 0,043) e com o de MCC (rho = 0,5515; p = 0,098), e correlação fraca com o de G-Mean (rho = 0,2970; p = 0,405), coeficientes estimados sobre dez modelos e, portanto, indicativos da tendência observada neste conjunto, e não estimativas populacionais precisas. A divergência é evidenciada de forma direta pelo próprio ranking, que produz três vencedores distintos entre os quatro critérios: Random Forest em ROC-AUC, HistGradient Boosting em G-Mean e Regressão Logística em MCC e Kappa. Esse resultado converge com a recomendação de De la Cruz Huayanay, Bazán e Russo (\citeyear{delacruz2025}) contra o uso isolado da ROC-AUC como critério de seleção de modelos em cenários de dados desbalanceados, e reforça a importância de portfólios de métricas complementares em estudos de churn, prática ainda pouco adotada em trabalhos aplicados ao dataset Olist.

A análise de importância de features revelou que frequência de compras, recência e diversidade de categorias (n\_categories) são os principais drivers preditivos do churn, respondendo por 47,26\% da importância total do modelo na execução de referência. A estabilidade da hierarquia entre partições sustenta esse achado: em trinta divisões estratificadas independentes, frequency ocupou a primeira posição em todas, n\_categories permaneceu entre as três primeiras em 29 e recency\_days em 22. O destaque de n\_categories, em empate técnico de importância com a recência, sugere que estratégias de cross-selling constituem mecanismo promissor de retenção proativa, resultado pouco explorado na literatura aplicada ao Olist e com implicação direta para o design de campanhas de marketing. A comparação dos rankings por valores SHAP entre os três modelos de melhor desempenho, apresentada na Tabela \ref{tab:shap-tres-modelos}, corrobora parcialmente essa hierarquia: a recência ocupa a segunda posição nos três estimadores e a frequência lidera nos dois modelos baseados em árvores, ao passo que a diversidade de categorias é confirmada entre as quatro primeiras por dois dos três modelos. Quanto à direção dos efeitos, os dependence plots apresentados na Seção \ref{subsec:perfil-retorna} evidenciam relações não lineares: a frequência e a diversidade de categorias passam a contribuir negativamente para o churn a partir de três pedidos e de três categorias distintas, a recência mantém relação monotônica crescente com o risco, e o parcelamento médio elevado associa-se a maior propensão ao abandono. Lida sob a ótica da retenção, essa configuração caracteriza o perfil do cliente que volta a comprar: maior frequência (2,56 contra 2,08 pedidos), recência baixa e maior diversidade de categorias (1,80 contra 1,54), caracterização que deve ser interpretada como hipótese a confirmar em bases de maior volume, dada a dimensão da classe Ativo.

A validação da representatividade do balanceamento via PSI e KS-test indicou que, apesar de o subconjunto retido corresponder a apenas 10,6\% da classe Churn completa (13,3\% da parcela presente no conjunto de treino), oito das nove features numéricas preservaram a distribuição estatística original, com desvio marginal restrito à variável avg\_installments (PSI = 0,1293, sem rejeição pelo teste KS). O resultado é compatível com a hipótese de que o modelo aprendeu padrões do Churn real, e não de uma versão sistematicamente distorcida da classe majoritária \cite{he2009}, ressalvado que a potência do teste KS é limitada pelo tamanho do subconjunto retido e que a verificação alcança o undersampling da classe majoritária, não a expansão sintética da classe minoritária. Essa etapa de validação, não localizada nos estudos de churn aplicados ao dataset Olist consultados neste trabalho, constitui contribuição metodológica adicional.

As principais contribuições do trabalho incluem: a comparação sistemática de dez algoritmos aplicados ao e-commerce brasileiro com protocolo multi-métrica; a validação de um pipeline anti-leakage com janela temporal dupla; a comparação empírica de quatro estratégias de tratamento do desbalanceamento, incluindo análise de robustez em trinta divisões independentes; a validação estatística da representatividade do subconjunto efetivamente balanceado; a identificação de n\_categories entre os três principais drivers preditivos do churn, pouco discutida na literatura aplicada ao Olist; a proposição de um critério de seleção de modelo baseado na captura em cobertura fixa, que neutraliza o tamanho do recorte e isola a qualidade da ordenação; e a documentação do efeito, sobre as estimativas de validação cruzada, da aplicação da reamostragem previamente ao particionamento, quantificado em queda de até 0,3251 na área sob a curva ROC para os algoritmos de maior capacidade.

O trabalho apresenta limitações decorrentes da filtragem de clientes com compra única, que reduz a base para 1.178 clientes, e da ausência de hiperparametrização exaustiva. A limitação mais relevante refere-se ao tamanho da classe minoritária: com aproximadamente cinquenta clientes Ativos na base e dez no conjunto de teste, as métricas dependentes de limiar apresentam elevada sensibilidade a variações amostrais, e cada cliente capturado na análise de cobertura fixa equivale a dez pontos percentuais de Recall, de modo que diferenças de uma única unidade entre modelos não são conclusivas. A análise de robustez conduzida na Seção \ref{sec:target-balanceamento} quantifica esse efeito: em trinta divisões estratificadas independentes, a ROC-AUC variou entre 0,500 e 0,822, com média de 0,6686, o que indica que o desempenho da execução de referência situa-se na faixa superior da distribuição. Cabe registrar que o comportamento das métricas de classificação sob amostras pequenas é lacuna explicitamente apontada por De la Cruz Huayanay, Bazán e Russo (\citeyear{delacruz2025}), cujo estudo empregou tamanhos de 5.000 e 10.000 observações, o que situa este trabalho fora do regime amostral em que aquelas recomendações foram estabelecidas. Registra-se ainda que nenhum dos estimadores de importância de features disponíveis resolve, isoladamente, a ordenação das variáveis nesse regime: a importância por permutação, calculada sobre um hold-out com dez instâncias da classe minoritária, apresenta desvios-padrão entre 44\% e 93\% das respectivas médias nas variáveis de média positiva, razão pela qual a estabilidade entre partições foi adotada como evidência. Adicionalmente, o dataset Olist cobre o período de 2016 a 2018, o que pode não refletir integralmente a dinâmica atual do mercado brasileiro de e-commerce. Por fim, registra-se que ordenações entre features de importância muito próxima podem variar marginalmente entre versões das bibliotecas utilizadas \cite{pedregosa2011}, ainda que com semente fixa, razão pela qual as conclusões deste trabalho são formuladas sobre os padrões robustos a essas variações.

Trabalhos futuros podem explorar: (1) feature engineering com dados de sessão e comportamento de navegação; (2) modelos de churn em tempo real com processamento via streaming; (3) hiperparametrização via otimização bayesiana (Optuna) \cite{imani2025}; (4) modelos de sobrevivência (Cox Proportional Hazards) para estimativa do tempo até o churn; (5) integração de SHAP por instância no contrato de serving da API para explicações personalizadas por cliente; (6) seleção de amostras típicas (clientes mais informativos) em substituição ao undersampling aleatório, conforme discutido na Seção \ref{subsec:clientes-informativos}; (7) calibração de probabilidades (Platt scaling ou regressão isotônica) para os modelos de boosting, visando estabilizar o ponto de corte; (8) exposição do modelo como ferramenta via protocolo MCP (Model Context Protocol), permitindo que analistas de negócio consultem scores de risco e explicações SHAP em linguagem natural, integrando o modelo a fluxos conversacionais de apoio à decisão; e (9) codificação alternativa das variáveis categóricas, por target encoding ou por agrupamento prévio de categorias de baixa frequência, em substituição ao Label Encoding adotado.
